% Generated from the audited Markdown source; responsible author: Carptopus.
\documentclass[11pt]{article}
\usepackage[a4paper,margin=28mm]{geometry}
\usepackage[T1]{fontenc}
\usepackage[utf8]{inputenc}
\usepackage{lmodern}
\usepackage{microtype}
\usepackage{amsmath,amssymb,amsthm,mathtools}
\usepackage{needspace}
\usepackage{xcolor}
\usepackage[colorlinks=true,linkcolor=blue!55!black,citecolor=blue!55!black,urlcolor=blue!65!black]{hyperref}
\usepackage{fancyhdr}

\newtheorem{theorem}{Theorem}[section]
\newtheorem{lemma}[theorem]{Lemma}
\newtheorem{corollary}[theorem]{Corollary}
\numberwithin{equation}{section}

\pagestyle{fancy}
\fancyhf{}
\fancyhead[L]{Carptopus}
\fancyhead[R]{Dual Fano partition degrees}
\fancyfoot[C]{\thepage}
\setlength{\headheight}{14pt}
\setlength{\parindent}{0pt}
\setlength{\parskip}{0.48em}
\setlength{\emergencystretch}{4em}
\urlstyle{same}

\title{The partition-representation degrees of the dual Fano matroid\\
and a six-symbol characterization of regularity}
\author{Carptopus\\\href{mailto:carptopus@163.com}{carptopus@163.com}}
\date{24 August 2026}

\hypersetup{
  pdftitle={The partition-representation degrees of the dual Fano matroid and a six-symbol characterization of regularity},
  pdfauthor={Carptopus},
  pdfsubject={A complete partition-degree classification for the dual Fano matroid and a six-symbol characterization of regular matroids},
  pdfkeywords={matroid, entropic matroid, partition representation, variable-strength orthogonal array, dual Fano matroid, regular matroid, orthogonal Latin square}
}

\begin{document}
\maketitle

\begin{center}
\fcolorbox{black!35}{black!4}{\parbox{0.88\textwidth}{\centering\small
Public Beta, version 0.1. The mathematical proof, complete manuscript, and
release artifacts have passed project-internal zero-trust audits;
external mathematical review and formal peer review are pending.}}
\end{center}

\begin{abstract}


For a finite matroid $M$, let $\chi(M)$ be the set of integers $v\ge 2$ for
which $M$ has a partition representation of degree $v$, equivalently a
variable-strength orthogonal array of level $v$. We prove that the dual Fano
matroid satisfies

\nopagebreak[4]
\[
\chi(F_7^*)=\{2^k:k\ge 1\}.
\]

The proof fixes one coordinate fibre, identifies the resulting contraction
with $M(K_4)$, applies Mat\'u\v{s}'s group normal form, and propagates that local
normal form through the remaining circuit equations. The propagation forces
every element of the underlying group to be an involution. Combining this
classification with the excluded-minor characterization of regular matroids
gives a fixed-alphabet characterization: a finite matroid is 6-entropic if
and only if it is regular. Moreover, among all integers $v\ge 2$, the
class of $v$-entropic matroids coincides with the class of regular matroids
if and only if $v=6$.

\end{abstract}

\section{Introduction}


A partition representation of a matroid is a nonlinear analogue of a linear
representation. It may be expressed equivalently as an almost-affine code, an
ideal secret-sharing representation, or a variable-strength orthogonal array
(VOA). If a matroid $M$ has rank function $r_M$, a degree-$v$
representation realizes the entropy vector

\nopagebreak[4]
\[
A\longmapsto r_M(A)\log v.
\]

Following Abbe and Spirkl, we call $M$ $v$-entropic when its rank function
is realized in this way \cite{abbe2019}. Following Chen, Cheng, and Bai, we write
$\chi(M)$ for the set of all admissible degrees \cite{chen2025}.

In the proof of Proposition 4.1, Mat\'u\v{s} showed that the Fano matroid has
precisely the power-of-two degrees \cite{matus1999}. For the dual Fano matroid, Abbe and
Spirkl gave an analytic proof that degree $3$ is impossible \cite[Lemma 9]{abbe2019};
Chen, Cheng, and Bai later obtained an independent finite-computation proof
\cite[Lemma 4]{chen2025}. The latter authors also posed the general question whether
$\chi(M)=\chi(M^*)$ for every matroid \cite[Conjecture 1]{chen2025}. Our first result
settles the complete degree set for this particular dual pair.

\Needspace{0.26\textheight}
\begin{theorem}
For the dual Fano matroid,

\nopagebreak[4]
\[
\chi(F_7^*)=\{2^k:k\ge 1\}.
\]

\end{theorem}

The argument is structural rather than enumerative. A fibre of a hypothetical
degree-$v$ representation of $F_7^*$ represents $M(K_4)$. Mat\'u\v{s}'s
classification of partition representations of $M(K_4)$ supplies a finite
group of order $v$. Three further circuit equations then force that group
to have exponent two.

The result has a consequence that is not restricted to a seven-element
matroid. Recall that a finite matroid is regular if it is representable over
every field, equivalently if it has no minor isomorphic to
$U_{2,4}$, $F_7$, or $F_7^*$.

\Needspace{0.26\textheight}
\begin{theorem}
A finite matroid is 6-entropic if and only if it is regular.

\end{theorem}

Chen, Cheng, and Bai proved that every regular matroid has a partition
representation of every degree \cite[Proposition 7]{chen2025}. The new direction in
Theorem 1.2 follows by excluding the three regular-matroid obstructions at
degree six: the obstruction $U_{2,4}$ is Euler's order-six exception for
orthogonal Latin squares, while $F_7$ and $F_7^*$ are excluded by their
power-of-two degree sets.

The alphabet size six is characterized by this property.

\Needspace{0.26\textheight}
\begin{corollary}
For an integer $v\ge 2$, the following are equivalent:

\begin{enumerate}
\item every $v$-entropic matroid is regular;
\item the $v$-entropic matroids are exactly the regular matroids;
\item $v=6$.

\end{enumerate}

\end{corollary}

Indeed, $F_7$ is a nonregular 2-entropic matroid. For every
$v\ge 3$, $v\ne 6$, the nonregular matroid $U_{2,4}$ is
$v$-entropic by the complete existence theorem for pairs of orthogonal
Latin squares.

The group slice used below is classical. Recent work of Kalantari and
Khazaei also uses Mat\'u\v{s}'s $M(K_4)$ normal form in its characterization of
4-entropic matroids \cite{kalantari2026}. Our contribution is the propagation across the
specific circuit system of $F_7^*$, the resulting complete degree
classification, and its degree-six consequences.

\section{Partition representations and circuit equations}


Let $M=(E,r)$ be a finite matroid of rank $r(E)$, and let $Q$ be a set
of size $v$. A variable-strength orthogonal array
$\mathrm{VOA}(M,v)$ is a $v^{r(E)}\times |E|$ array $T$ over
$Q$ such that, for each $A\subseteq E$, every row appearing in the
projection $T(A)$ appears exactly $v^{r(E)-r(A)}$ times. Its coordinate
level partitions form a partition representation of degree $v$, and the
two notions are equivalent \cite{matus1999,chen2025}. The entropy characterization of Chen,
Cheng, and Bai shows that $M$ is $v$-entropic exactly when such an array
exists \cite[Theorem 1 and the discussion following it]{chen2025}; see also \cite[Proposition 2.3]{kalantari2026}.

We use three standard consequences.

\Needspace{0.26\textheight}
\begin{lemma}[basis coordinates]
If $B$ is a basis of $M$, every tuple
in $Q^B$ occurs exactly once in $T(B)$. Hence each coordinate outside
$B$ is a function of the basis coordinates.

\end{lemma}

\begin{proof}
Since $r(B)=|B|=r(E)$, the VOA multiplicity on $B$ is one.
There are $v^{r(E)}=|Q^B|$ rows, so the projection is a bijection.
\end{proof}


\Needspace{0.26\textheight}
\begin{lemma}[circuit quasigroups]
Let $C$ be a circuit of $M$ and
$e\in C$. In every degree-$v$ representation, the coordinate $X_e$
is a quasigroup function of the coordinates $X_{C\setminus\{e\}}$.

\end{lemma}

\begin{proof}
The set $C\setminus\{e\}$ is independent and
$r(C)=|C|-1$. Every tuple on $C\setminus\{e\}$ therefore has a unique
extension to the joint support on $C$. Repeating this with each coordinate
singled out gives the quasigroup property.
\end{proof}


Here a $k$-ary quasigroup is a map $q:Q^k\to Q$ for which fixing any
(k-1) arguments makes the remaining one a bijection.

\Needspace{0.26\textheight}
\begin{lemma}[minor closure]
If $v\in\chi(M)$ and $N$ is a minor of
$M$, then $v\in\chi(N)$.

\end{lemma}

\begin{proof}
Deletion is obtained by projecting the array and deduplicating
rows. Contraction by an element is obtained by fixing one coordinate value
and restricting to that fibre. Iterating these two operations gives every
minor. This is \cite[Theorem 2]{chen2025}.
\end{proof}


Coordinate relabellings preserve a VOA. More precisely, one may independently
apply a bijection $Q\to Q$ to each coordinate. In Mat\'u\v{s}'s terminology this
is partition isotopy. The same relabellings transform all circuit equations
simultaneously; they cannot be chosen separately for different fibres.

\section{A group slice inside \texorpdfstring{$F_7^*$}{F7*}}


We use the labelling of $F_7^*$ from \cite{chen2025}. Its seven circuit-hyperplanes are

\nopagebreak[4]
\[
1235,\qquad 1246,\qquad 1347,\qquad 1567,\qquad
2367,\qquad 2457,\qquad 3456.             \tag{3.1}
\]

The set (1234) is a basis. By Lemmas 2.1 and 2.2, a degree-$v$
representation can be written in basis coordinates as

\nopagebreak[4]
\[
\begin{aligned}
X_1&=a,&X_2&=b,&X_3&=c,&X_4&=d,\\
X_5&=F(a,b,c),&
X_6&=G(a,b,d),&
X_7&=H(a,c,d),
\end{aligned}                         \tag{3.2}
\]

where $F,G,H$ are ternary quasigroups on an alphabet $Q$ of size $v$.

Fix an arbitrary value $a=a_0$. Restricting to this fibre is the VOA
contraction by element (1). The circuit-hyperplanes through (1) give the
four triangles

\nopagebreak[4]
\[
235,\qquad 246,\qquad 347,\qquad 567.  \tag{3.3}
\]

They are precisely the four triangles of a labelled copy of $M(K_4)$.
For example, a binary representation of $F_7^*$ may be chosen with

\nopagebreak[4]
\[
5=e_1+e_2+e_3,\qquad
6=e_1+e_2+e_4,\qquad
7=e_1+e_3+e_4.
\]

After contraction by $e_1$, the six nonzero columns are

\nopagebreak[4]
\[
e_2,e_3,e_4,e_2+e_3,e_2+e_4,e_3+e_4,
\]

so there are no loops, parallel elements, or additional triangles.

Mat\'u\v{s}'s classification now supplies the key local normal form.

\Needspace{0.26\textheight}
\begin{lemma}[group normal form]
After fixed relabellings of the six
coordinate alphabets $2,\ldots,7$, there is a finite group
$(\Gamma,\cdot)$ of order $v$ such that, on the fibre $a=a_0$,

\nopagebreak[4]
\[
F(a_0,b,c)=b^{-1}c,\qquad
G(a_0,b,d)=b^{-1}d,\qquad
H(a_0,c,d)=c^{-1}d.                  \tag{3.4}
\]

The group is not assumed commutative.

\end{lemma}

\begin{proof}
Mat\'u\v{s}'s Proposition 3.1 states that every partition representation
of $M(K_4)$ is partition-isotopic to the canonical representation obtained
from a finite group of the same order \cite{matus1999}. One orientation of that canonical
model is

\nopagebreak[4]
\[
(x,y,z,xy^{-1},yz^{-1},zx^{-1}).     \tag{3.5}
\]

Put $b=x^{-1}$, $c=y^{-1}$, and $d=z^{-1}$. Keep the output
coordinates corresponding to $xy^{-1}$ and $yz^{-1}$, and invert the
coordinate corresponding to $zx^{-1}$. After matching these coordinates
with the triangles in (3.3), the three functions become exactly (3.4).
These substitutions are six fixed coordinate bijections. We apply them to
the entire representation and use the identity relabelling on coordinate
1; no relabelling depends on $a$.
\end{proof}


\section{Propagation through the remaining circuits}


The three circuit-hyperplanes $2367,2457,3456$ propagate the normalized
fibre to every value of $a$. The point requiring care is that a circuit
equation is initially known on the joint support of its four coordinates.
In each of the following steps, the normalized fibre covers the complete
three-dimensional input domain of the relevant circuit quasigroup.

\Needspace{0.26\textheight}
\begin{lemma}
There is, for each $a\in Q$, a permutation
$T_a:\Gamma\to\Gamma$ such that

\nopagebreak[4]
\[
G(a,b,d)=b^{-1}T_a(d),\qquad
H(a,c,d)=c^{-1}T_a(d).               \tag{4.1}
\]

\end{lemma}

\begin{proof}
Orient the circuit $2367$ toward coordinate $7$. On the
normalized fibre, write $g=b^{-1}d$, so $d=bg$. Then

\nopagebreak[4]
\[
h=c^{-1}d=c^{-1}bg.
\]

As $b,c,g$ independently range over $\Gamma$, this determines the entire
ternary circuit quasigroup. Therefore, for every basis tuple,

\nopagebreak[4]
\[
H(a,c,d)=c^{-1}bG(a,b,d).            \tag{4.2}
\]

The left-hand side is independent of $b$. Thus, for fixed $a,d$, the
product $bG(a,b,d)$ is independent of $b$; call it $T_a(d)$. Since
$G$ is a quasigroup in $d$, $T_a$ is a permutation. Equation (4.1)
follows.
\end{proof}


\Needspace{0.26\textheight}
\begin{lemma}
For every $a\in Q$, there is an element
$s_a\in\Gamma$ such that

\nopagebreak[4]
\[
\begin{aligned}
F(a,b,c)&=b^{-1}s_a^{-1}c,\\
G(a,b,d)&=b^{-1}s_a d,\\
H(a,c,d)&=c^{-1}s_a d.
\end{aligned}               \tag{4.3}
\]

\end{lemma}

\begin{proof}
Orient $2457$ toward coordinate $7$. On the normalized fibre,
put $f=b^{-1}c$, so $c=bf$. Then

\nopagebreak[4]
\[
h=c^{-1}d=f^{-1}b^{-1}d.
\]

The variables $b,d,f$ independently range over $\Gamma$, so the whole
circuit quasigroup is determined and

\nopagebreak[4]
\[
H(a,c,d)=F(a,b,c)^{-1}b^{-1}d.       \tag{4.4}
\]

Substituting (4.1) into (4.4) gives

\nopagebreak[4]
\[
F(a,b,c)^{-1}=c^{-1}T_a(d)d^{-1}b.
\]

The left-hand side is independent of $d$. Hence $T_a(d)d^{-1}$ is
constant in $d$; denote it by $s_a$. Thus $T_a(d)=s_ad$, and (4.3)
follows.
\end{proof}


\Needspace{0.26\textheight}
\begin{lemma}
Every element $s_a$ in (4.3) is an involution.

\end{lemma}

\begin{proof}
Orient $3456$ toward coordinate $6$. On the normalized fibre,
$f=b^{-1}c$ gives $b=cf^{-1}$, and hence

\nopagebreak[4]
\[
g=b^{-1}d=fc^{-1}d.
\]

The variables $c,d,f$ independently range over $\Gamma$. Consequently
the complete circuit equation is

\nopagebreak[4]
\[
G(a,b,d)=F(a,b,c)c^{-1}d.            \tag{4.5}
\]

Using (4.3) on the two sides of (4.5) yields

\nopagebreak[4]
\[
b^{-1}s_ad=b^{-1}s_a^{-1}d.
\]

Thus $s_a=s_a^{-1}$, or $s_a^2=e$.
\end{proof}


\Needspace{0.26\textheight}
We can now prove the first main theorem.

\begin{proof}[Proof of Theorem 1.1]
Fix $b,c\in\Gamma$. Since $F$ is a ternary
quasigroup, the map

\nopagebreak[4]
\[
a\longmapsto F(a,b,c)=b^{-1}s_a^{-1}c
\]

is a bijection. Hence $a\mapsto s_a$ is a bijection from the alphabet
$Q$ onto $\Gamma$. By Lemma 4.3, every element $x\in\Gamma$ satisfies
$x^2=e$. Therefore

\nopagebreak[4]
\[
xy=(xy)^{-1}=y^{-1}x^{-1}=yx
\]

for all $x,y\in\Gamma$. The group is elementary abelian of order $2^k$
for some $k\ge1$, proving that every degree is a power of two.

Conversely, $F_7^*$ is binary. A binary representing matrix has the same
column ranks over every extension field $\mathbb F_{2^k}$, and the associated
linear array is a degree-$2^k$ partition representation. Hence every
$2^k$ occurs.
\end{proof}


\section{Six-symbol entropy and regularity}


\Needspace{0.26\textheight}
We now prove the global consequence.

\begin{proof}[Proof of Theorem 1.2]
Suppose first that $M$ is 6-entropic. By the
entropy--VOA--partition equivalence, $6\in\chi(M)$. Lemma 2.3 implies that
every minor of $M$ also has degree six.

The matroid $U_{2,4}$ has a degree-$v$ representation exactly when there
is an index-one orthogonal array $\mathrm{OA}(2,4,v)$, equivalently a
pair of orthogonal Latin squares of order $v$. Such a pair does not exist
for $v=6$; equivalently, $6\notin\chi(U_{2,4})$ \cite[Example 4]{chen2025}. The proof
of Mat\'u\v{s}'s Proposition 4.1 gives $\chi(F_7)=\{2^k:k\ge1\}$ \cite{matus1999}, and
Theorem 1.1 gives
the same degree set for $F_7^*$. Thus none of

\nopagebreak[4]
\[
U_{2,4},\qquad F_7,\qquad F_7^*
\]

is a minor of $M$. Tutte's excluded-minor theorem for regular matroids
implies that $M$ is regular \cite[Theorem 6.6.6]{oxley2011}.

Conversely, let $M$ be regular. A totally unimodular representation of
$M$, reduced modulo $v$, defines a VOA over the ring $\mathbb Z_v$ for
every integer $v\ge2$: every basis determinant is $\pm1$, hence remains
a unit. This is proved explicitly in \cite[Lemma 3 and Proposition 7]{chen2025}. In
particular $6\in\chi(M)$, so $M$ is 6-entropic.
\end{proof}


\begin{proof}[Proof of Corollary 1.3]
The implication (2)$\Rightarrow$(1) is immediate,
and Theorem 1.2 gives (3)$\Rightarrow$(2).

It remains to show that (1) forces $v=6$. If $v=2$, the nonregular Fano
matroid $F_7$ is $v$-entropic. If $v\ge3$ and $v\ne6$, there exists a
pair of orthogonal Latin squares of order $v$, and hence the nonregular
matroid $U_{2,4}$ is $v$-entropic. The existence statement is the theorem
of Bose, Shrikhande, and Parker together with the classical small orders
\cite[Example 4]{chen2025}; \cite{bose1960}. Therefore (1) fails for every $v\ne6$.
\end{proof}


\section{Scope and relation to previous work}


The proof uses several established ingredients whose roles should be kept
separate from the new argument.

\begin{enumerate}
\item Mat\'u\v{s} introduced partition representations, classified the
$M(K_4)$ representations by finite groups, and determined
$\chi(F_7)$ \cite{matus1999}.
\item Abbe and Spirkl gave an analytic proof that
$3\notin\chi(F_7^*)$ \cite[Lemma 9]{abbe2019}. Chen, Cheng, and Bai developed the
VOA and entropy formulation, proved minor closure and arbitrary-degree
representability of regular matroids, and later gave an independent
finite-computation proof of the same degree-three obstruction \cite{chen2025}.
\item Kalantari and Khazaei recently used the same $M(K_4)$ group slice in
their degree-four characterization \cite{kalantari2026}. Thus the group normal form itself
is not claimed as new.
\item Tutte's excluded-minor characterization and the orthogonal-Latin-square
existence theorem are classical \cite{oxley2011,bose1960}.

\end{enumerate}

The new point is that the three remaining circuit equations of $F_7^*$
make the local group normal form global and force exponent two. This yields
the complete degree set of $F_7^*$, rather than the exclusion of one fixed
degree. Theorems 1.2 and Corollary 1.3 are consequences at the level of all
finite matroids.

This result verifies
$\chi(F_7)=\chi(F_7^*)$ for the Fano pair, but it does not prove the general
duality conjecture of \cite{chen2025}. Nor does it classify partition representations up
to partition isotopy; it classifies only the possible degrees.

The literature comparison made for this manuscript found no direct statement
of Theorem 1.1, Theorem 1.2, or Corollary 1.3 in the sources above. This is a
bounded prior-art statement, not a claim of absolute global priority.

\section{Computational calibration}


The proof is noncomputational. A small verification script accompanies the
research record only to detect multiplication-order errors in the propagation
formulas. It checks the full $F_7^*$ VOA conditions for elementary abelian
2-groups of orders $2,4,8$, and checks that cyclic groups of orders $4,6$
and the nonabelian group $S_3$ fail at the final circuit as predicted.
Independent audit calculations also used the dihedral and quaternion groups
of order eight as noncommutative negative controls.

These finite checks do not establish any universal assertion and are not used
in the proof.

\section*{Acknowledgement of AI assistance}


The author used Codex as an AI-assisted research and writing tool for
literature organization, symbolic checking, adversarial proof review, and
manuscript preparation. The author selected the claims, reviewed the evidence,
and assumes responsibility for the mathematical content and the final text.

\begin{thebibliography}{9}
\footnotesize
\raggedright

\bibitem{matus1999}
F.~Mat\'u\v{s},
\emph{Matroid representations by partitions},
Discrete Math. 203 (1999), 169--194.
\href{https://doi.org/10.1016/S0012-365X(99)00004-7}{doi:10.1016/S0012-365X(99)00004-7}.

\bibitem{chen2025}
Q.~Chen, M.~Cheng, and B.~Bai,
\emph{Matroidal entropy functions: constructions, characterizations, and representations},
IEEE Trans. Inf. Theory 71(4) (2025), 3237--3249.
\href{https://doi.org/10.1109/TIT.2024.3355942}{doi:10.1109/TIT.2024.3355942}.

\bibitem{abbe2019}
E.~Abbe and S.~Spirkl,
\emph{Entropic matroids and their representation},
Entropy 21(10) (2019), 948.
\href{https://doi.org/10.3390/e21100948}{doi:10.3390/e21100948}.

\bibitem{kalantari2026}
M.~H. Kalantari and S.~Khazaei,
\emph{Four-entropic matroids are quaternary},
arXiv:2608.20553v1 (2026).
\href{https://arxiv.org/abs/2608.20553}{arXiv:2608.20553}.

\bibitem{oxley2011}
J.~Oxley,
\emph{Matroid Theory}, second edition,
Oxford University Press, 2011.

\bibitem{bose1960}
R.~C. Bose, S.~S. Shrikhande, and E.~T. Parker,
\emph{Further results on the construction of mutually orthogonal Latin squares and the falsity of Euler's conjecture},
Canad. J. Math. 12 (1960), 189--203.
\href{https://doi.org/10.4153/CJM-1960-016-5}{doi:10.4153/CJM-1960-016-5}.

\end{thebibliography}

\end{document}
